\begin{figure}[htbp]
	\centering
	\resizebox{\textwidth}{!}{%
		\begin{tikzpicture}[
			asta/.style={line width=1.5pt, line cap=round,
				line join=round},
			nodo/.style={circle, fill=white, draw=black,
				inner sep=0pt, minimum size=6pt,
				line width=1.2pt},
			forza/.style={->, >=stealth, line width=1.0pt,
				gray!60!black},
			]
			
			%============================================================
			% PANNELLO A — schema di carico (sinistra)
			%============================================================
			\begin{scope}[shift={(0,0)}, scale=0.62]
				
				\coordinate (A) at (0,0);
				\coordinate (B) at (8,0);
				\coordinate (C) at (4,3);
				\coordinate (D) at (4,0);
				
				%% carico trasversale p (perpendicolare ad AC, sopra)
				\draw[gray!60!black, thin]
				($(A)+0.9*(-0.6,0.8)$)--($(C)+0.9*(-0.6,0.8)$);
				\draw[gray!60!black, thin]
				($(A)+0.25*(-0.6,0.8)$)--($(C)+0.25*(-0.6,0.8)$);
				\foreach \t in {0.0, 0.2, 0.4, 0.6, 0.8, 1.0}{
					\draw[forza]
					($(A)!{\t}!(C)+0.9*(-0.6,0.8)$)
					--($(A)!{\t}!(C)+0.25*(-0.6,0.8)$);
				}
				\node[above left, font=\small, gray!60!black]
				at ($(A)!0.5!(C)+1.0*(-0.6,0.8)$){$p$};
				
				%% carico longitudinale q (radente sotto AC, da A verso C)
				\foreach \t in {0.05, 0.28, 0.51, 0.74}{
					\draw[forza]
					($(A)!{\t}!(C)+0.3*(0.6,-0.8)$)
					--($(A)!{\t+0.18}!(C)+0.3*(0.6,-0.8)$);
				}
				\node[below right, font=\small, gray!60!black]
				at ($(A)!0.5!(C)+0.4*(0.6,-0.8)$){$q$};
				
				%% aste
				\draw[asta] (A)--(C);
				\draw[asta] (C)--(B);
				\draw[asta] (A)--(D);
				\draw[asta] (D)--(B);
				\draw[asta] (C)--(D);
				
				%% cerniera A
				\draw[line width=1.2pt]
				(A)--++(-0.4,-0.65)--++(0.8,0)--cycle;
				\draw[line width=1.2pt](-0.55,-0.65)--(0.55,-0.65);
				\foreach \x in {-0.45,-0.25,-0.05,0.15,0.35}
				\draw[thin](\x,-0.65)--++(-.18,-.22);
				
				%% carrello B
				\draw[line width=1.2pt]
				(B)--++(-0.4,-0.50)--++(0.8,0)--cycle;
				\draw[line width=1.2pt]
				($(B)+(-0.50,-0.50)$)--($(B)+(0.50,-0.50)$);
				\draw[line width=1.2pt]($(B)+(-0.25,-0.65)$) circle(3.5pt);
				\draw[line width=1.2pt]($(B)+( 0.25,-0.65)$) circle(3.5pt);
				\draw[line width=1.2pt]
				($(B)+(-0.55,-0.82)$)--($(B)+(0.55,-0.82)$);
				\foreach \x in {-0.45,-0.25,-0.05,0.15,0.35}
				\draw[thin]($(B)+(\x,-0.82)$)--++(-.18,-.22);
				
				%% nodi
				\node[nodo] at (A){};  \node[nodo] at (B){};
				\node[nodo] at (C){};  \node[nodo] at (D){};
				
				%% etichette nodi
				\node[left,  font=\normalsize] at (A){$A$};
				\node[right, font=\normalsize] at (B){$B$};
				\node[right, font=\normalsize] at (C){$C$};
				\node[below, font=\normalsize] at (D){$D$};
				
				%% etichetta pannello
				\node[font=\small] at (4,-1.5) {$(a)$};
				
			\end{scope}
			
			%============================================================
			% PANNELLO B — asta isolata (destra)
			%============================================================
			\begin{scope}[shift={(7,0)}, scale=0.62]
				
				%% asta isolata orizzontale nel rif. locale
				\draw[asta] (0,0)--(8,0);
				
				%% tratteggio sotto l'asta (riferimento locale)
				\draw[thin, dashed] (0,-0.3)--(8,-0.3);
				
				%% carico trasversale p
				\draw[gray!60!black, thin] (0,2.5)--(8,2.5);
				\draw[black, thin] (0,1.85)--(8,1.85);
				\foreach \x in {0,1,2,3,4,5,6,7,8}{
					\draw[forza] (\x,2.5)--(\x,1.85);
				}
				\node[above, font=\small, gray!60!black] at (4,2.5){$p$};
				
				%% carico longitudinale q (radente, da A verso C)
				\foreach \t in {0.5, 2.0, 3.5, 5.0, 6.5}{
					\draw[forza] (\t,0.4)--(\t+0.8,0.4);
				}
				\node[above, font=\small, gray!60!black] at (4,0.4){$q$};
				
				%% sforzo normale Ne (esce da C verso destra)
				\draw[->, >=stealth, line width=1.5pt, black]
				(8,0)--(9.5,0);
				\node[above, font=\normalsize] at (9.5,0){$N^e$};
				
				%% sforzo normale N(0) (esce da A verso sinistra)
				\draw[->, >=stealth, line width=1.5pt, black]
				(0,0)--(-1.5,0);
				\node[above, font=\normalsize] at (-1.5,0){$N(0)$};
				
				%% reazioni trasversali (coda in A e C, verso l'alto)
				\draw[->, >=stealth, gray!60!black, line width=1.0pt]
				(0,0)--(0,1.35);
				\node[left, font=\small] at (0,1.2){$R_A^\perp$};
				\draw[->, >=stealth, gray!60!black, line width=1.0pt]
				(8,0)--(8,1.35);
				\node[right, font=\small] at (8,1.2){$R_C^\perp$};
				
				%% nodi
				\node[nodo] at (0,0){};
				\node[nodo] at (8,0){};
				
				%% etichette nodi
				\node[below left,  font=\normalsize] at (0,0){$A$};
				\node[below right, font=\normalsize] at (8,0){$C$};
				
				%% etichetta pannello
				\node[font=\small] at (4,-1.5) {$(b)$};
				
			\end{scope}
			
		\end{tikzpicture}
	}% fine resizebox
	\caption{Asta reticolare $AC$ soggetta a carichi generici.
		\emph{(a)}: schema strutturale della capriata con
		carico trasversale $p$ (perpendicolare ad $AC$,
		sopra l'asta) e carico longitudinale $q$
		(radente ad $AC$ da $A$ verso $C$).
		\emph{(b)}: asta $AC$ isolata nel riferimento locale
		con $z^e$ diretto da $A$ verso $C$;
		$N^e$ e $N(0)$ sono gli sforzi normali uscenti
		dalle due estremità; $R_A^\perp$ e $R_C^\perp$
		sono le reazioni trasversali alle cerniere.}
	\label{fig:asta_caricata}
\end{figure}
